\section{Introduction}\label{sec:intro}

Formalizing mathematics has been slow. A non-trivial theorem could take
dozens or hundreds of person-hours to formalize, and a large one years
of work by a team. This has changed. In September 2026 a language model
wrote a proof of Fermat's Last Theorem in Lean, 13 million lines in 11
days~\cite{anthropic2026flt}. Writing the proof is no longer the
bottleneck.

What remains is the statement. A formal proof proves what its statement
says, and checking that the statement, and the definitions it depends
on, say what was meant needs a person. Generated formalizations are not
written to be read, so this check is still slow. As
\citet{hales2026reliability} puts it, a proof in Lean should not be
accepted until a human has checked statement fidelity, a task much
easier than checking the proof. For Fermat's Last Theorem the check was
easy, because the generated statement could be compared mechanically
with one that people had already written and checked in Mathlib. Most
results have no such statement to compare with. We would like the
check to stay easy anyway.

In this paper we take one step towards this, in one area of quantum
computing where formalization is wanted: completeness of equational
theories for circuits. A quantum circuit denotes a unitary matrix. A
set of equations between circuits is \emph{sound} if both sides of
each equation denote the same matrix, and \emph{complete} if every
equation between circuits with the same matrix follows from it. For
example, Selinger's fifteen relations for the gates $\omega$, $H$, $S$
and $CZ$ are complete for $n$-qubit Clifford circuits, for every
$n$~\cite{selinger2015clifford}. Soundness is a finite check.
Completeness says that the equations present the group or monoid of
matrices that the circuits generate, and it is much harder to prove.
It is also what a compiler wants: a complete set of rules decides the
equality of circuits, and an optimizer built on it misses no
equation.

There are a few standard ways to prove it. One defines a normal form,
shows that every circuit can be rewritten to one, and shows that
different normal forms have different matrices. Another starts from an
exact synthesis algorithm, which turns a matrix into a word, and shows
by induction on a measure of the matrix that every generator applied to
it can be absorbed into the algorithm's output. A third moves a known
complete theory along a translation between circuit languages. In each
case the core of the proof is a case analysis, and the papers say so
in their page counts: the extended version of
Hadamard-$\Pi$~\cite{fang2026hadamard} has 116 pages, Makary, Ross and
Selinger support their real stabilizer theorem with a supplement of 76
pages of derivations~\cite{makary2021real}, and Cl\'ement's theorem
for real Clifford+$CH$ circuits has a 13-page text and an appendix of
264 pages~\cite{clement2026realcliffordch}. Such appendices are seldom
checked in full by a conference referee, and an author cannot easily
check them again after a change either.

These papers have one thing in common that helps: they share a
semantics. A circuit is a word over gates, it denotes a matrix, and the
theorem says that the relations present the group, or monoid, that the
matrices of the circuits form. So the form of the statement is the same
in every paper, and it can be fixed once, by code that people write and
check once. We call such code a scaffold. Our earlier Agda
library~\cite{anonymous2027popl} is the syntactic half of one: words,
circuits on $n$ wires with their structural rules, normal forms, and
what it means for a relation set to present a group or a subgroup, with
any monoid or group as the semantics. Its examples stopped where the
matrices begin. For the Clifford+$T$ monoids it showed only that the
gate-set presentation is isomorphic to an amalgamated product of
smaller presentations, and it said what was missing: rings and unitary
groups that the standard library does not have, and then the
completeness proof of each paper. This paper supplies that half. A
layer of matrix semantics gives, over the rings of exact synthesis, the
unitary and orthogonal groups of the two-level papers and the operators
that circuits denote, and a second addition proves rules independent.
Every completeness theorem below then has one form with two
parameters: a relation set over the library's words, which for
circuits is its circuit syntax, and a group of matrices, which comes
from the layer, or, for the qudit Clifford circuits, an algebraic model
of it. A person who knows the form reads every statement the same way,
and checks per paper only what the scaffold cannot fix: the relations
as transcribed, the matrices of the gates, and the one line that states
the theorem. For the real stabilizer development these come to under
500 of its 12{,}010 lines. Agda checks the rest.

The scaffold also fixes what the author of a new theory has to do,
which is the other half of the cost. The author writes the same two
parameters, and then the mathematics of the paper's own proof along
one of four routes: a normal form with its rewriting and its
uniqueness argument, an exact synthesis algorithm with its Main Lemma,
a decomposition by width, or an encoding into a theory already known
to be complete; a group built from presented pieces is assembled by
the library's products and extensions instead. Much of the rest is
done by machine: soundness is decided by evaluation, the structural
rules are sound once for every gate set, the induction of the
synthesis route and the argument of the encoding route are modules
written once, and the derivations and case analyses that a paper
gives as computations are generated by scripts and checked by Agda.
What a paper gives as an idea is still written. For the real
stabilizer development half of the lines are generated and a fifth are
the paper's mathematics (Section~\ref{sec:example}).

We did this for fifteen completeness proofs, written between 2014 and
2026, and for the minimality theorem of a sixteenth paper.
Table~\ref{tab:corpus} lists them. Each development states its result
against a semantics: the relations present the group of matrices that
the generators generate, or they present a submonoid of a matrix group,
at every width or dimension. The one exception is the qudit Clifford
group of \citet{bian2026qupit}, whose relations are shown to present an
algebraic model of that group (Section~\ref{sec:qupit}). A development
has about 26{,}000 lines on average, and all of them together about
1{,}500 modules and 420{,}000 lines, on a library of 17{,}500 lines,
about a third of it added for this work. They typecheck in safe
Agda~2.8 with no postulates, in about 40 minutes
(Table~\ref{tab:costs}), and one module states the sixteen theorems.

Safe Agda is also constructive: no excluded middle, no choice, no
extensionality, and no real numbers in the standard library. Lean's
Mathlib, on which the proof of Fermat's Last Theorem is built, uses
choice and extensionality throughout; our proofs assume nothing of the
kind. That makes them harder to write, and it makes each of them a
program (Section~\ref{sec:why-agda}).

The most important finding is about Hadamard-$\Pi$, a POPL~2026 paper
on equational reasoning for quantum programs~\cite{fang2026hadamard}.
Its Theorem~4.14 says that the relations of its Figure~6 present the
group $O_n(\Z[\frac1{\sqrt2}])$ of orthogonal matrices over
$\Z[\frac1{\sqrt2}]$. The proof is incomplete: two of its lemmas are
false, one case of its Main Lemma relies on them, and that case fails
on an explicit state in dimension~9 (Section~\ref{sec:gap}). The
theorem is still true. An unpublished thesis of Cl\'ement from
2019~\cite{clement2019thesis} presents the same group by other
relations and handles the case by other diagrams; we proved the two
relation sets equivalent, formalized his argument for the missing
case, and later his whole proof. Cl\'ement's 2026 paper on real
Clifford+$CH$ circuits imports Theorem~4.14 as its Theorem~4.4, so its
main theorem rested on this gap too; our formalization of that paper
discharges the import. One more statement is false as printed: the
relations of Amy, Ross and Wesley's theorem on Toffoli--Hadamard
circuits~\cite{amy2024toffoli} omit one that their proof uses, and
without it they are not complete. The other papers are correct in what
we formalized, up to misprints, some of them in statements that proofs
use (Section~\ref{sec:findings}, Appendix~\ref{sec:appendix-errata}).

The qudit Clifford development was written mainly by a person, over
three months. LLM coding agents wrote the other fifteen, and the
scripts that generated their long derivations, within four weeks, most
of them in less than a day each, because almost every claim in these
papers can be checked by a computation before it is formalized, and
Agda checks the result after (Section~\ref{sec:method}). We think that
scaffolds made by people, on which the models build, are one way of
working with language models on mathematics: the people check the
scaffold once and a short statement per result, they supply the
mathematics that a route asks for, and the machine checks the rest.

\paragraph{Contributions.}
\begin{itemize}
\item An extension of the library of~\citet{anonymous2027popl} from
  abstract semantics to matrices: a layer of matrix semantics over the
  rings of exact synthesis, and a framework for the independence of
  rules. With it every completeness or minimality theorem has the form
  of Section~\ref{sec:theorem-form}, and a reader checks a short
  statement per paper (Section~\ref{sec:scaffold}).
\item Machine-checked proofs of fifteen completeness theorems for
  quantum circuits and two-level matrix groups, and of one minimality
  theorem, each stated against a matrix semantics or, for qudit
  Clifford circuits, an algebraic model (Section~\ref{sec:corpus}).
\item A gap in the completeness proof of Hadamard-$\Pi$, with explicit
  counterexamples to two of its lemmas and to one case of its Main
  Lemma, and two machine-checked repairs (Section~\ref{sec:gap}).
\item A list of the other errors that the formalization found, in the
  papers and in our own transcriptions (Section~\ref{sec:findings},
  Appendix~\ref{sec:appendix-errata}).
\item The proof routes behind these results as they appear in Agda,
  what each asks of an author, and what is generated and checked by
  machine, with one development followed from its gate set to its
  theorem (Section~\ref{sec:architecture}).
\item The methods we used to produce the proofs, and what they cost in
  Agda, including the performance problems we met and how we solved
  them (Sections~\ref{sec:method} and~\ref{sec:engineering}).
\end{itemize}

Section~\ref{sec:background} recalls presentations and the two main
kinds of completeness proof; Section~\ref{sec:conclusion} says what is
not formalized.
